% !TEX root = ../proyect.tex

\chapter{Conceptos iniciales}\label{cap:conceptos}

En este capítulo se presentan los conceptos fundamentales necesarios para comprender el desarrollo del proyecto LinceUS Assistant.

\section{Procesamiento de Lenguaje Natural (NLP)}\label{sec:nlp}

% TODO: Explicar qué es NLP
% - Definición
% - Aplicaciones
% - Técnicas principales utilizadas en el proyecto

\section{Sistemas conversacionales}\label{sec:sistemas-conversacionales}

% TODO: Explicar qué son los asistentes conversacionales
% - Chatbots vs asistentes inteligentes
% - Arquitectura típica
% - Componentes principales

\section{Aprendizaje automático aplicado a diálogos}\label{sec:ml-dialogos}

% TODO: Conceptos de ML en sistemas de diálogo
% - Clasificación de intenciones
% - Extracción de entidades
% - Políticas de diálogo

\section{Modelos de lenguaje grandes (LLMs)}\label{sec:llms}

% TODO: Explicar qué son los LLMs
% - GPT, Llama, Mistral
% - Text-to-SQL
% - Generación de lenguaje natural

\section{Bases de datos relacionales}\label{sec:bd-relacionales}

% TODO: Conceptos de BD relacionales
% - Modelo entidad-relación
% - Normalización
% - Consultas SQL

\section{Búsqueda semántica y RAG}\label{sec:busqueda-semantica}

% TODO: Explicar búsqueda semántica
% - Embeddings y vectores
% - Similitud coseno
% - RAG (Retrieval-Augmented Generation)
% - pgvector
